\documentclass[10pt,a4paper]{article}
\usepackage[T1]{fontenc}
\usepackage[margin=2.2cm]{geometry}
\usepackage{amsmath,amssymb,amsthm}
\usepackage{booktabs}
\usepackage{array}
\usepackage{xurl}
\usepackage[hidelinks]{hyperref}
\newcommand{\FNfoneZero}{0.309}
\newcommand{\FNfoneOne}{0.672}
\newcommand{\FNfoneOAOne}{0.408}
\newcommand{\FNconfOne}{0.302}
\newcommand{\phatThree}{0.384}
\newcommand{\PowESVZero}{0.280}
\newcommand{\PowESVHalf}{0.141}
\newcommand{\PowESVWorst}{0.052}
\newcommand{\PowOAHalf}{0.204}
\newcommand{\PowOAWorst}{0.150}
\newcommand{\ExtCalls}{30\%}
\newcommand{\Nquestions}{1,109}
\newcommand{\NholdClosed}{188}
\newcommand{\SEhold}{0.022}
\newcommand{\HoldNse}{1.2}
\newcommand{\Nscores}{twelve}
\newcommand{\NconfViolOne}{ten}
\newcommand{\ConfAllBy}{0.3}
\newcommand{\MaxInsitu}{0.101}
\newcommand{\MaxTransfer}{0.100}
\newcommand{\FNbgeZero}{0.236}
\newcommand{\FNbgeOne}{0.731}
\newcommand{\PowBGEZero}{0.443}
\newcommand{\PowBGEHalf}{0.173}
\newcommand{\PowBGEOne}{0.097}
\newcommand{\RobPMEsvIn}{0.099}
\newcommand{\RobPMEsvTr}{0.094}
\newcommand{\RobPMAllIn}{0.101}
\newcommand{\RobPMAllTr}{0.100}
\newcommand{\RobPMConf}{0.302}
\newcommand{\RobPTEsvIn}{0.099}
\newcommand{\RobPTEsvTr}{0.094}
\newcommand{\RobPTAllIn}{0.101}
\newcommand{\RobPTAllTr}{0.104}
\newcommand{\RobPTConf}{0.277}
\newcommand{\RobHMEsvIn}{0.117}
\newcommand{\RobHMEsvTr}{0.104}
\newcommand{\RobHMAllIn}{0.126}
\newcommand{\RobHMAllTr}{0.126}
\newcommand{\RobHMConf}{0.299}
\newcommand{\RobHTEsvIn}{0.115}
\newcommand{\RobHTEsvTr}{0.108}
\newcommand{\RobHTAllIn}{0.125}
\newcommand{\RobHTAllTr}{0.125}
\newcommand{\RobHTConf}{0.279}
\newcommand{\ClimN}{1,381}
\newcommand{\ClimFoneMax}{0.60}
\newcommand{\ClimFoneZeroMin}{0.29}
\newcommand{\ClimFoneZeroMax}{0.42}
\newcommand{\ClimConfMax}{0.240}
\newcommand{\ClimCertMax}{0.103}
\newcommand{\ClimDensePowZero}{0.256}
\newcommand{\ClimDensePowOne}{0.195}
\newcommand{\RetroY}{2010}
\newcommand{\RetroYears}{2008, 2010 and 2012}
\newcommand{\RetroNclosed}{310}
\newcommand{\RetroNfilled}{326}
\newcommand{\RetroNnei}{416}
\newcommand{\RetroVisible}{56\%}
\newcommand{\RetroBestPow}{68\%}
\newcommand{\RetroBestPrec}{96\%}
\newcommand{\RetroMaxFN}{0.102}
\newcommand{\RetroFoneMin}{28\%}
\newcommand{\RetroFoneMax}{40\%}
\newcommand{\RetroNdossier}{139}
\newcommand{\RetroLag}{4}
\newcommand{\LLMfone}{0.629}
\newcommand{\LLMfn}{0.378}
\newcommand{\LLMaucFull}{0.687}
\newcommand{\LLMaucDel}{0.477}
\newcommand{\LLMfoneOne}{0.617}
\newcommand{\LLMpowZero}{0.191}

\newcommand{\SuppNdos}{139}
\newcommand{\SuppLagMin}{1}
\newcommand{\SuppLagMax}{8}
\newcommand{\SuppTauDos}{0.030}
\newcommand{\SuppClimN}{1,381}
\newcommand{\SuppClimCal}{690}
\newcommand{\SuppClimEval}{691}
\newcommand{\SuppR}{500}
\newcommand{\SuppM}{100}
\newcommand{\SuppDelta}{0.05}
\newcommand{\SuppCVsel}{0.718}


\newtheorem{theorem}{Theorem}
\newtheorem{corollary}{Corollary}
\newtheorem{proposition}{Proposition}
\theoremstyle{definition}
\newtheorem{definition}{Definition}
\theoremstyle{remark}
\newtheorem*{remark}{Remark}

\newcommand{\open}{\textup{\textsc{open}}}
\newcommand{\closed}{\textup{\textsc{closed}}}
\renewcommand{\thetable}{S\arabic{table}}
\renewcommand{\thesection}{S\arabic{section}}

\title{Supplementary Material for\\``Absence of Evidence Is Not a Research Gap:\\Certified Gap Claims under Corpus Incompleteness''}
\author{Anh Hoa Le, Ly Van Khoa Phan, Dinh Thanh Nguyen,\\Duc Hoang Nguyen, Dinh Anh Ho, Long Truong\\FPT University, Vietnam}
\date{}

\begin{document}
\sloppy
\maketitle

\noindent This document supports the main paper. Section~\ref{sec:proofs} gives complete statements and proofs, including the proof of the Corollary, which the paper states without proof. Section~\ref{sec:impl} lists implementation details needed to reproduce the experiments. Sections~\ref{sec:unc} to~\ref{sec:retro} report the full results behind the summary tables of the paper: every score, rule, deletion rate and evaluation setting. All tables are generated by \texttt{make\_supp.py} from the same result files as the tables of the paper, so the numbers agree with it. The code is available at \url{https://github.com/Khoa-Hoa-Technology-Solution-Company/research-paper-gap} (directory \texttt{ESV-Gap/gapclose}).

Throughout, FN is the false novelty rate (share of \closed{} questions declared \open{}), power is the share of \open{} questions declared \open{}, $\alpha=0.1$ is the target, and values above $\alpha+0.005$ are printed in bold.

\section{Proofs}\label{sec:proofs}

\subsection{Setting}
A question $q$ is \closed{} if the full literature $\mathcal{C}$ contains a witness document for it and \open{} otherwise. Under the independent witness-deletion model, each witness document of $\mathcal{C}$ is absent from the local corpus $\mathcal{C}_p$ independently with probability $p$, and other documents are unaffected. A closure score $s(q;\mathcal{K})$ is computed against a corpus $\mathcal{K}$; the rule declares \open{} iff $s(q;\mathcal{C}_p)<\tau$. A score is monotone if $\mathcal{K}'\subseteq\mathcal{K}$ implies $s(q;\mathcal{K}')\le s(q;\mathcal{K})$. Given calibration scores $S_1,\dots,S_n$ of \closed{} questions, sorted increasingly, the conformal threshold is $\tau=S_{(k)}$ with $k=\lfloor\alpha(n+1)\rfloor$, and $\tau=-\infty$ if $k=0$.

\begin{proposition}[Absence is not identifiable]
Let $q$ be a question for which $\mathcal{C}_p$ contains no witness, and let $d^*$ be a document that would be a witness for $q$. If $0<p<1$, the full corpora $\mathcal{C}_{\mathrm{open}}=\mathcal{C}_p$ and $\mathcal{C}_{\mathrm{closed}}=\mathcal{C}_p\cup\{d^*\}$ both produce the observation $\mathcal{C}_p$ with positive probability. Any rule that is a function of $\mathcal{C}_p$ alone therefore gives the same decision in both worlds.
\end{proposition}
\begin{proof}
In $\mathcal{C}_{\mathrm{open}}$, $q$ has no witness and $\mathcal{C}_p=\mathcal{C}_{\mathrm{open}}$ with probability one, since only witnesses are deleted. In $\mathcal{C}_{\mathrm{closed}}$, $d^*$ is deleted with probability $p>0$, and every other witness in it is kept with probability $1-p>0$, independently, so $\mathcal{C}_p$ is observed with positive probability. A deterministic rule maps the same observation to the same decision. If the rule declares \open{} on $\mathcal{C}_p$, it commits false novelty with positive probability in the closed world; if it declares \closed{}, it is wrong in the open world. Bounding false novelty therefore requires information on $p$ or from outside $\mathcal{C}_p$.
\end{proof}

\begin{theorem}[In-situ calibration]
Suppose the calibration questions are scored against the same incomplete corpus $\mathcal{C}_p$ as the test question, their \closed{} labels come from outside $\mathcal{C}_p$, and the calibration and test questions are exchangeable. Then $P(s(q_{\mathrm{test}};\mathcal{C}_p)<\tau\mid q_{\mathrm{test}}\ \closed)\le\alpha$ for every $p$, known or not.
\end{theorem}
\begin{proof}
Deletion acts on documents and does not depend on question identity, so the joint law of the $n+1$ pairs (question, state of its witnesses in $\mathcal{C}_p$) is invariant under permutations of the questions. The $n+1$ scores are therefore exchangeable, and the rank of the test score among them is uniform on $\{1,\dots,n+1\}$ (ties broken at random). The event $S_{\mathrm{test}}<S_{(k)}$ requires the test score to have rank at most $k$, so its probability is at most $k/(n+1)\le\alpha$. Labels from outside $\mathcal{C}_p$ are needed because a label inferred from $\mathcal{C}_p$ would select calibration questions whose witnesses survived, which breaks exchangeability with the test question.
\end{proof}

\begin{theorem}[Transfer calibration]
Assume that the calibration and test questions are exchangeable before deletion, that the calibration scores are computed on a complete benchmark after deleting each of its witnesses independently with probability $\hat p$, and that the target corpus follows the deletion model with rate $p\le\hat p$. If $s$ is monotone, then the test false novelty on $\mathcal{C}_p$ is at most $\alpha$.
\end{theorem}
\begin{proof}
Construct $\mathcal{C}_{\hat p}$ from $\mathcal{C}_p$ by deleting each surviving witness independently with probability $r=(\hat p-p)/(1-p)$ (for $p=1$ set $\mathcal{C}_{\hat p}=\mathcal{C}_p$). A witness survives both steps with probability $(1-p)(1-r)=1-\hat p$, independently across witnesses, so $\mathcal{C}_{\hat p}$ has the law of the deletion model at rate $\hat p$, and $\mathcal{C}_{\hat p}\subseteq\mathcal{C}_p$ by construction. Monotonicity gives $s(q;\mathcal{C}_{\hat p})\le s(q;\mathcal{C}_p)$ pointwise, hence
$P(s(q;\mathcal{C}_p)<\tau)\le P(s(q;\mathcal{C}_{\hat p})<\tau)$.
At rate $\hat p$, the test score and the calibration scores are produced by the same deletion process applied to exchangeable questions, so the argument of Theorem~1 bounds the right-hand side by $\alpha$.
\end{proof}

\begin{corollary}[Estimated coverage]
Let $m$ reference documents be witnesses of known \closed{} questions, distinct from the witnesses of the screened questions, and let $x$ of them be missing from $\mathcal{C}_p$. Let $\hat p$ be the one-sided $(1-\delta)$ Clopper-Pearson upper bound for $x/m$, and compute the transfer threshold with this $\hat p$. Under the conditions of Theorem~2, the false novelty is at most $\alpha+\delta$.
\end{corollary}
\begin{proof}
Under the deletion model the $m$ reference documents are missing independently with probability $p$, so $x\sim\mathrm{Bin}(m,p)$ and the Clopper-Pearson bound satisfies $P(\hat p<p)\le\delta$. Because the reference documents are distinct from the witnesses of the test question, their deletion is independent of the state of those witnesses, so $\hat p$ is independent of $s(q_{\mathrm{test}};\mathcal{C}_p)$. The calibration scores are drawn afresh at rate $\hat p$ from the complete benchmark. Conditioning on $\hat p$,
\[
\mathrm{FN}=E\big[P(s(q_{\mathrm{test}};\mathcal{C}_p)<\tau_{\hat p}\mid\hat p)\big]
\le \alpha\,P(\hat p\ge p)+1\cdot P(\hat p<p)\le\alpha+\delta,
\]
where the first bound applies Theorem~2 for each fixed value $\hat p\ge p$.
\end{proof}

\begin{remark}
Two details of our simulation differ slightly from these assumptions. First, the $m=\SuppM{}$ reference documents are drawn with replacement from the witnesses of the calibration questions, so a document can be counted twice, and a document can be a witness of both a calibration and a test question. Each draw is still missing with probability $p$, so the estimate is unbiased, but the draws are not fully independent and the Clopper-Pearson bound is then approximate. Second, under \emph{targeted} deletion the deletion probability of a witness is $\min\{1,p\,w_d\}$, where $w_d$ is proportional to the strongest ESV score of $d$ for any of its questions plus $0.05$, normalized to mean one; this violates the independence across witness strength on purpose. Table~\ref{tab:phat} shows that the estimate stays above the realized deletion rate on average at every $p$.
\end{remark}

\begin{remark}[Which scores are monotone]
A score that takes a maximum over all available documents of a fixed pool can only decrease when documents are removed. This covers CoMention, BM25 top-1, Dense top-1, BGE scores, the LLM judge and ESV-Scope over the full retrieval pool, and also the combination $\max\{s_{\mathrm{loc}},s_{\mathrm{ext}}\}$ because $s_{\mathrm{ext}}$ does not depend on the local corpus. Scores computed over the top $k$ documents, or by a learned head over features of the top $k$, are not monotone in general: removing a document lets another enter the top $k$. In the training split, under two nested random deletion sets, NLI verification over the top five documents increased for 18 of 809 questions, ESV-Scope over the top five for 15 and ESV-Learned for 46, while the max-over-pool scores never increased.
\end{remark}

\section{Implementation details}\label{sec:impl}

\paragraph{Benchmarks.} GapClose-SciFact reads each SciFact claim as the question whether published evidence on it exists. Claims with at least one \textsc{support} or \textsc{contradict} rationale are \closed{}, with their rationale documents as witnesses; \textsc{nei} claims are \open{}. The official training split (809 questions) is used for development and the development split (300 questions) as test set. GapClose-ClimateFEVER (\SuppClimN{} questions) treats each of the five evidence sentences of a claim as a document titled by its Wikipedia article; claims labeled \textsc{supports} or \textsc{refutes} are \closed{}, \textsc{not enough info} claims are \open{}, and the 154 \textsc{disputed} claims are dropped. Resampled splits use \SuppClimCal{} calibration and \SuppClimEval{} evaluation questions.

\paragraph{Retrieval.} For each question the pool $P(q)$ is the top ten documents of reciprocal rank fusion ($k_{\mathrm{RRF}}=60$) over BM25 and the bi-encoder \texttt{sentence-transformers/all-MiniLM-L6-v2}, each contributing its top 20. BM25 and Dense baselines use the top-1 score of their own ranking.

\paragraph{Verifier and scope.} Sentences are scored by \url{MoritzLaurer/DeBERTa-v3-large-mnli-fever-anli-ling-wanli}, sentence as premise and question as hypothesis. Content terms are lower-cased tokens longer than two characters, outside a stop list, with light deterministic suffix stripping (for example \emph{-ation}, \emph{-ing}, \emph{-es}); numbers are matched by the pattern \verb|\d+(\.\d+)?|. The ESV-Scope exponent and depth were selected on the training split by five-fold cross-validation (Table~\ref{tab:cv}); ESV-Learned is a logistic regression over eight verifier and scope features of the top three documents (selected CV macro-F1 \SuppCVsel{}).

\paragraph{Other baselines.} BGE dense uses \texttt{bge-large-en-v1.5} (top-1 cosine over the corpus), the reranker \texttt{bge-reranker-base} (maximum over the pool). The LLM judge runs \texttt{Qwen2.5-1.5B-Instruct} (ONNX, fp16) with the system message ``You are a careful scientific fact checker.'' and the user message ``Text: \{abstract\}. Does this text contain evidence that supports or refutes the following claim? Claim: \{claim\}. Answer with Yes or No.''; the score of a document is the probability mass on \emph{Yes} relative to \emph{Yes} plus \emph{No}, and the question score is the maximum over the pool. Abstracts are truncated to 320 tokens.

\paragraph{External evidence.} The OpenAlex \texttt{works} endpoint is queried with the stop-word-filtered question text (at most 400 characters), requesting ten works; works without an abstract are dropped, abstracts are split into at most 30 sentences longer than 20 characters, and the same verifier and scope factor are applied. Failed or rate-limited requests are retried and never cached as empty results. In the retrospective experiment, works published after the cutoff are discarded.

\paragraph{Certification protocol.} Each setting uses $R=\SuppR{}$ repetitions with a fixed random seed. In each repetition a single global deletion pattern is drawn, scores of calibration and evaluation \closed{} questions are computed on it, and five thresholds are derived from the calibration half: the F1-optimal threshold on the complete corpus, split conformal on the complete corpus, in-situ conformal on the incomplete corpus, transfer conformal at the estimated $\hat p$ ($m=\SuppM{}$, $\delta=\SuppDelta{}$), and transfer at $\hat p=1$. \open{} questions have no witnesses and are unaffected by deletion.

\paragraph{Compute.} All models run through ONNX Runtime with the published weights (DirectML on an RTX 3060 laptop GPU). Retrieval and NLI scoring of SciFact take about one hour, the BGE baselines about 15 minutes, the LLM judge about 75 minutes, and the four certification settings about 50 minutes in parallel. External retrieval needed about 1,100 OpenAlex searches.

\begin{table}[h]
\centering
\caption{Five-fold cross-validation on the SciFact training split (macro-F1, mean $\pm$ SD over folds). $a$: scope exponent, $k$: number of verified documents. The test split was not used.}
\label{tab:cv}
\setlength{\tabcolsep}{3pt}
\begin{tabular}{lccc}
\toprule
Variant & $k=3$ & $k=5$ & $k=10$ \\
\midrule
ESV-Scope, $a=0$ & 0.651 $\pm$ 0.029 & 0.631 $\pm$ 0.044 & 0.615 $\pm$ 0.055 \\
ESV-Scope, $a=0.25$ & 0.685 $\pm$ 0.016 & 0.687 $\pm$ 0.009 & 0.682 $\pm$ 0.013 \\
ESV-Scope, $a=0.5$ & 0.694 $\pm$ 0.009 & 0.682 $\pm$ 0.016 & 0.696 $\pm$ 0.016 \\
ESV-Scope, $a=1$ & 0.684 $\pm$ 0.020 & 0.685 $\pm$ 0.015 & 0.704 $\pm$ 0.023 \\
ESV-Learned & 0.718 $\pm$ 0.019 & 0.704 $\pm$ 0.017 & 0.698 $\pm$ 0.016 \\
\midrule
CoMention & \multicolumn{3}{c}{0.660 $\pm$ 0.026} \\
BM25 & \multicolumn{3}{c}{0.643 $\pm$ 0.013} \\
Dense (MiniLM) & \multicolumn{3}{c}{0.654 $\pm$ 0.040} \\
NLI-verify ($k=5$) & \multicolumn{3}{c}{0.631 $\pm$ 0.044} \\
\bottomrule
\end{tabular}

\end{table}

\section{Uncertified accuracy}\label{sec:unc}

The paper reports two passes over the test split. Table~\ref{tab:runs} gives both. Run 1 used DeBERTa-v3-small and untuned ESV defaults; run 2 used the large verifier and the configuration selected in Table~\ref{tab:cv}. Run 1 is reported for disclosure. After it we switched to the large verifier and selected every ESV setting by cross-validation on the training split only; no setting was chosen on test results. Table~\ref{tab:paired} gives the paired bootstrap comparison of ESV-Learned with the baselines of run 2: no difference is significant. Table~\ref{tab:unc} adds metrics not shown in the paper.

\begin{table}[h]
\centering
\caption{The two passes over the SciFact test split (300 questions).}
\label{tab:runs}
\setlength{\tabcolsep}{3pt}
\begin{tabular}{>{\raggedright\arraybackslash}p{4.1cm}cccc}
\toprule
 & \multicolumn{2}{c}{Run 1 (small)} & \multicolumn{2}{c}{Run 2 (large)} \\
\cmidrule(lr){2-3}\cmidrule(lr){4-5}
Method & F1 & FN & F1 & FN \\
\midrule
CoMention & 0.669 & 0.197 & 0.669 & 0.197 \\
BM25 & 0.660 & 0.197 & 0.660 & 0.197 \\
Dense (MiniLM) & 0.631 & 0.319 & 0.631 & 0.319 \\
NLI-verify & 0.533 & 0.346 & 0.663 & 0.239 \\
ESV (run 1: Scope, $a=0.5$, $k=5$; run 2: Learned, $k=3$) & 0.648 & 0.324 & 0.687 & 0.234 \\
\bottomrule
\end{tabular}

\end{table}

\begin{table}[h]
\centering
\caption{Paired comparison on the test split, ESV-Learned minus baseline. $\Delta$F1 with paired bootstrap 95\% interval; McNemar counts of questions only ESV-Learned got right ($b$) and only the baseline got right ($c$).}
\label{tab:paired}
\setlength{\tabcolsep}{3pt}
\begin{tabular}{lccc}
\toprule
Baseline & $\Delta$F1 & 95\% CI & McNemar $b$ / $c$ \\
\midrule
CoMention & +0.018 & [-0.042, +0.077] & 38 / 36 \\
BM25 & +0.027 & [-0.045, +0.096] & 54 / 50 \\
Dense (MiniLM) & +0.056 & [-0.013, +0.122] & 63 / 45 \\
NLI-verify & +0.024 & [-0.041, +0.086] & 43 / 37 \\
\bottomrule
\end{tabular}

\end{table}

\begin{table}[h]
\centering
\caption{Further uncertified metrics on the test split. Wit.\ doc: share of quoted witness documents that are gold witnesses; Gr.\ rec.: share of \closed{} questions closed with a gold witness document; FC hard: share of hard-negative documents of \open{} questions that were used to close the question.}
\label{tab:unc}
\setlength{\tabcolsep}{2.5pt}
\begin{tabular}{lcccccc}
\toprule
Method & Acc. & FN 95\% CI & FC 95\% CI & Wit.\ doc & Gr.\ rec. & FC hard \\
\midrule
CoMention & 0.700 & [0.14, 0.26] & [0.38, 0.56] & 0.596 & 0.479 & 0.283 \\
BM25 & 0.693 & [0.14, 0.25] & [0.40, 0.58] & 0.768 & 0.617 & 0.345 \\
Dense (MiniLM) & 0.647 & [0.25, 0.38] & [0.32, 0.50] & 0.727 & 0.495 & 0.348 \\
NLI-verify & 0.687 & [0.18, 0.30] & [0.35, 0.53] & 0.671 & 0.511 & 0.184 \\
ESV-Learned & 0.707 & [0.18, 0.30] & [0.31, 0.48] & 0.792 & 0.606 & 0.182 \\
\bottomrule
\end{tabular}

\end{table}

\clearpage
\section{Certification under incompleteness, all scores}\label{sec:cert}

Tables~\ref{tab:fnf1} to~\ref{tab:fntr} give the mean false novelty over \SuppR{} resampled splits under random deletion, for every score and deletion rate. The paper's Table~V is the ESV-Scope row of these tables. Table~\ref{tab:q95} shows that the guarantee is on the mean: in an individual split, the realized false novelty of the certified rules can exceed $\alpha$, as expected for a finite evaluation set.

\begin{table}[h]
\centering
\caption{Mean false novelty, F1-optimal threshold, by deletion rate $p$.}
\label{tab:fnf1}
\setlength{\tabcolsep}{2.5pt}
\begin{tabular}{lccccccc}
\toprule
Score & $0$ & $0.1$ & $0.2$ & $0.3$ & $0.5$ & $0.7$ & $1$ \\
\midrule
CoMention & \textbf{0.259} & \textbf{0.293} & \textbf{0.321} & \textbf{0.353} & \textbf{0.417} & \textbf{0.486} & \textbf{0.580} \\
BM25 & \textbf{0.229} & \textbf{0.261} & \textbf{0.289} & \textbf{0.330} & \textbf{0.397} & \textbf{0.466} & \textbf{0.580} \\
Dense & \textbf{0.308} & \textbf{0.342} & \textbf{0.370} & \textbf{0.402} & \textbf{0.465} & \textbf{0.530} & \textbf{0.626} \\
BGE dense & \textbf{0.212} & \textbf{0.255} & \textbf{0.298} & \textbf{0.349} & \textbf{0.432} & \textbf{0.520} & \textbf{0.656} \\
BGE rerank & \textbf{0.236} & \textbf{0.286} & \textbf{0.334} & \textbf{0.381} & \textbf{0.475} & \textbf{0.575} & \textbf{0.731} \\
LLM judge & \textbf{0.309} & \textbf{0.338} & \textbf{0.367} & \textbf{0.396} & \textbf{0.458} & \textbf{0.514} & \textbf{0.617} \\
NLI-verify & \textbf{0.256} & \textbf{0.273} & \textbf{0.302} & \textbf{0.332} & \textbf{0.391} & \textbf{0.452} & \textbf{0.536} \\
ESV-Scope & \textbf{0.309} & \textbf{0.343} & \textbf{0.380} & \textbf{0.415} & \textbf{0.489} & \textbf{0.564} & \textbf{0.672} \\
\midrule
CoMention +OA & \textbf{0.267} & \textbf{0.290} & \textbf{0.318} & \textbf{0.341} & \textbf{0.391} & \textbf{0.444} & \textbf{0.516} \\
Dense +OA & \textbf{0.239} & \textbf{0.260} & \textbf{0.284} & \textbf{0.307} & \textbf{0.354} & \textbf{0.403} & \textbf{0.478} \\
BGE rerank +OA & \textbf{0.244} & \textbf{0.283} & \textbf{0.323} & \textbf{0.361} & \textbf{0.438} & \textbf{0.518} & \textbf{0.646} \\
ESV-Scope +OA & \textbf{0.257} & \textbf{0.276} & \textbf{0.288} & \textbf{0.307} & \textbf{0.335} & \textbf{0.369} & \textbf{0.408} \\
\bottomrule
\end{tabular}

\end{table}

\begin{table}[h]
\centering
\caption{Mean false novelty, split conformal on the complete corpus.}
\label{tab:fnconf}
\setlength{\tabcolsep}{2.5pt}
\begin{tabular}{lccccccc}
\toprule
Score & $0$ & $0.1$ & $0.2$ & $0.3$ & $0.5$ & $0.7$ & $1$ \\
\midrule
CoMention & 0.068 & 0.087 & \textbf{0.108} & \textbf{0.128} & \textbf{0.171} & \textbf{0.213} & \textbf{0.279} \\
BM25 & 0.098 & \textbf{0.122} & \textbf{0.144} & \textbf{0.165} & \textbf{0.212} & \textbf{0.254} & \textbf{0.325} \\
Dense & 0.098 & \textbf{0.117} & \textbf{0.132} & \textbf{0.151} & \textbf{0.185} & \textbf{0.223} & \textbf{0.281} \\
BGE dense & 0.098 & \textbf{0.136} & \textbf{0.171} & \textbf{0.208} & \textbf{0.282} & \textbf{0.359} & \textbf{0.480} \\
BGE rerank & 0.097 & \textbf{0.137} & \textbf{0.174} & \textbf{0.212} & \textbf{0.287} & \textbf{0.363} & \textbf{0.486} \\
LLM judge & 0.099 & \textbf{0.119} & \textbf{0.134} & \textbf{0.151} & \textbf{0.189} & \textbf{0.226} & \textbf{0.286} \\
NLI-verify & 0.099 & \textbf{0.114} & \textbf{0.131} & \textbf{0.145} & \textbf{0.181} & \textbf{0.216} & \textbf{0.274} \\
ESV-Scope & 0.099 & \textbf{0.117} & \textbf{0.135} & \textbf{0.153} & \textbf{0.196} & \textbf{0.236} & \textbf{0.302} \\
\midrule
CoMention +OA & 0.062 & 0.078 & 0.096 & \textbf{0.112} & \textbf{0.147} & \textbf{0.182} & \textbf{0.235} \\
Dense +OA & 0.097 & \textbf{0.116} & \textbf{0.129} & \textbf{0.146} & \textbf{0.177} & \textbf{0.209} & \textbf{0.258} \\
BGE rerank +OA & 0.098 & \textbf{0.126} & \textbf{0.153} & \textbf{0.179} & \textbf{0.233} & \textbf{0.287} & \textbf{0.368} \\
ESV-Scope +OA & 0.098 & \textbf{0.107} & \textbf{0.119} & \textbf{0.127} & \textbf{0.150} & \textbf{0.168} & \textbf{0.200} \\
\bottomrule
\end{tabular}

\end{table}

\begin{table}[h]
\centering
\caption{Mean false novelty, in-situ conformal (Theorem~1).}
\label{tab:fnin}
\setlength{\tabcolsep}{2.5pt}
\begin{tabular}{lccccccc}
\toprule
Score & $0$ & $0.1$ & $0.2$ & $0.3$ & $0.5$ & $0.7$ & $1$ \\
\midrule
CoMention & 0.068 & 0.085 & 0.092 & 0.088 & 0.084 & 0.077 & 0.090 \\
BM25 & 0.098 & 0.101 & 0.100 & 0.099 & 0.098 & 0.098 & 0.098 \\
Dense & 0.098 & 0.100 & 0.098 & 0.098 & 0.100 & 0.098 & 0.097 \\
BGE dense & 0.098 & 0.101 & 0.098 & 0.099 & 0.099 & 0.098 & 0.098 \\
BGE rerank & 0.097 & 0.101 & 0.099 & 0.098 & 0.100 & 0.097 & 0.098 \\
LLM judge & 0.099 & 0.101 & 0.098 & 0.098 & 0.098 & 0.099 & 0.099 \\
NLI-verify & 0.099 & 0.099 & 0.099 & 0.098 & 0.100 & 0.101 & 0.099 \\
ESV-Scope & 0.099 & 0.098 & 0.099 & 0.098 & 0.098 & 0.099 & 0.071 \\
\midrule
CoMention +OA & 0.062 & 0.078 & 0.091 & 0.093 & 0.089 & 0.087 & 0.084 \\
Dense +OA & 0.097 & 0.101 & 0.097 & 0.098 & 0.098 & 0.099 & 0.097 \\
BGE rerank +OA & 0.098 & 0.100 & 0.099 & 0.099 & 0.099 & 0.099 & 0.099 \\
ESV-Scope +OA & 0.098 & 0.098 & 0.099 & 0.099 & 0.100 & 0.100 & 0.100 \\
\bottomrule
\end{tabular}

\end{table}

\begin{table}[h]
\centering
\caption{Mean false novelty, transfer with estimated $\hat p$ (Theorem~2).}
\label{tab:fntr}
\setlength{\tabcolsep}{2.5pt}
\begin{tabular}{lccccccc}
\toprule
Score & $0$ & $0.1$ & $0.2$ & $0.3$ & $0.5$ & $0.7$ & $1$ \\
\midrule
CoMention & 0.068 & 0.078 & 0.079 & 0.077 & 0.072 & 0.074 & 0.090 \\
BM25 & 0.091 & 0.090 & 0.088 & 0.088 & 0.089 & 0.090 & 0.098 \\
Dense & 0.093 & 0.091 & 0.088 & 0.090 & 0.093 & 0.093 & 0.097 \\
BGE dense & 0.088 & 0.085 & 0.083 & 0.085 & 0.089 & 0.090 & 0.098 \\
BGE rerank & 0.088 & 0.082 & 0.082 & 0.081 & 0.087 & 0.088 & 0.098 \\
LLM judge & 0.094 & 0.092 & 0.087 & 0.087 & 0.088 & 0.091 & 0.099 \\
NLI-verify & 0.094 & 0.090 & 0.090 & 0.088 & 0.091 & 0.093 & 0.099 \\
ESV-Scope & 0.094 & 0.089 & 0.087 & 0.087 & 0.090 & 0.090 & 0.071 \\
\midrule
CoMention +OA & 0.062 & 0.075 & 0.081 & 0.081 & 0.080 & 0.077 & 0.084 \\
Dense +OA & 0.093 & 0.092 & 0.089 & 0.089 & 0.091 & 0.093 & 0.097 \\
BGE rerank +OA & 0.091 & 0.087 & 0.086 & 0.086 & 0.089 & 0.092 & 0.099 \\
ESV-Scope +OA & 0.095 & 0.092 & 0.093 & 0.093 & 0.095 & 0.096 & 0.100 \\
\bottomrule
\end{tabular}

\end{table}

\begin{table}[h]
\centering
\caption{ESV-Scope, mean / 95th percentile of false novelty over the \SuppR{} splits.}
\label{tab:q95}
\setlength{\tabcolsep}{3pt}
\begin{tabular}{lcccc}
\toprule
Rule & $p=0$ & $p=0.3$ & $p=0.5$ & $p=1$ \\
\midrule
In-situ & 0.099 / 0.140 & 0.098 / 0.138 & 0.098 / 0.135 & 0.071 / 0.144 \\
Transfer & 0.094 / 0.133 & 0.087 / 0.128 & 0.090 / 0.130 & 0.071 / 0.144 \\
\bottomrule
\end{tabular}

\end{table}

Table~\ref{tab:power} gives the power of transfer calibration for every score and rate, and the worst-case rule $\hat p=1$. Table~\ref{tab:phat} reports the coverage estimate: the mean $\hat p$ lies above the realized deletion rate at every $p<1$, by about 0.03 to 0.09, which is the margin that the Clopper-Pearson bound adds with \SuppM{} reference documents. The last column is the share of evaluation questions for which ESV-Scope +OA had to query OpenAlex, that is, whose local score fell below the threshold.

\begin{table}[h]
\centering
\caption{Power of transfer calibration (estimated $\hat p$) by deletion rate, and power of the worst-case rule at $p=0.5$.}
\label{tab:power}
\setlength{\tabcolsep}{2.5pt}
\begin{tabular}{lcccccccc}
\toprule
Score & $0$ & $0.1$ & $0.2$ & $0.3$ & $0.5$ & $0.7$ & $1$ & $\hat p{=}1$ \\
\midrule
CoMention & 0.300 & 0.264 & 0.216 & 0.184 & 0.143 & 0.127 & 0.114 & 0.114 \\
BM25 & 0.299 & 0.251 & 0.214 & 0.183 & 0.137 & 0.111 & 0.089 & 0.088 \\
Dense & 0.282 & 0.241 & 0.213 & 0.201 & 0.181 & 0.171 & 0.162 & 0.162 \\
BGE dense & 0.397 & 0.321 & 0.266 & 0.229 & 0.174 & 0.131 & 0.102 & 0.102 \\
BGE rerank & 0.443 & 0.345 & 0.280 & 0.235 & 0.173 & 0.129 & 0.097 & 0.097 \\
LLM judge & 0.191 & 0.162 & 0.140 & 0.130 & 0.114 & 0.103 & 0.086 & 0.084 \\
NLI-verify & 0.249 & 0.201 & 0.174 & 0.149 & 0.122 & 0.102 & 0.081 & 0.081 \\
ESV-Scope & 0.280 & 0.238 & 0.207 & 0.180 & 0.141 & 0.110 & 0.057 & 0.052 \\
\midrule
CoMention +OA & 0.288 & 0.278 & 0.244 & 0.210 & 0.161 & 0.132 & 0.115 & 0.115 \\
Dense +OA & 0.299 & 0.264 & 0.231 & 0.209 & 0.171 & 0.156 & 0.143 & 0.144 \\
BGE rerank +OA & 0.440 & 0.360 & 0.301 & 0.266 & 0.214 & 0.179 & 0.145 & 0.146 \\
ESV-Scope +OA & 0.305 & 0.275 & 0.254 & 0.235 & 0.204 & 0.179 & 0.151 & 0.150 \\
\bottomrule
\end{tabular}

\end{table}

\begin{table}[h]
\centering
\caption{Coverage estimate and cost, ESV-Scope, resampled splits, random deletion.}
\label{tab:phat}
\setlength{\tabcolsep}{3pt}
\begin{tabular}{ccccc}
\toprule
$p$ & realized rate & mean $\hat p$ & SD of power & ext.\ queries \\
\midrule
$0$ & 0.000 & 0.030 & 0.037 & 31\% \\
$0.1$ & 0.100 & 0.161 & 0.036 & 30\% \\
$0.2$ & 0.200 & 0.280 & 0.034 & 30\% \\
$0.3$ & 0.302 & 0.384 & 0.033 & 30\% \\
$0.5$ & 0.500 & 0.586 & 0.031 & 30\% \\
$0.7$ & 0.699 & 0.772 & 0.033 & 30\% \\
$1$ & 1.000 & 1.000 & 0.047 & 31\% \\
\bottomrule
\end{tabular}

\end{table}

\section{Robustness of the certified rules}\label{sec:robust}

Table~\ref{tab:robust} expands the paper's Table~IV to every score: the largest mean false novelty over the seven deletion rates, for both certified rules and all four evaluation settings. Under resampling, which matches the exchangeability assumption, all values stay within Monte Carlo error of the target, also under targeted deletion. In the train-to-test setting many scores exceed the target, including at $p=0$ where nothing is deleted (Tables~\ref{tab:fnin} and~\ref{tab:fntr} are the resampled counterpart). The evaluation set there contains only \NholdClosed{} \closed{} questions, so the binomial standard error of a single false novelty estimate is about \SEhold{}; we read the excess as a possible shift between the two SciFact splits that this experiment cannot separate from noise.

\begin{table}[h]
\centering
\caption{Largest mean false novelty over all deletion rates. Res.: resampled splits; T$\to$T: calibrate on the training split, evaluate on the test split.}
\label{tab:robust}
\setlength{\tabcolsep}{2pt}
\begin{tabular}{lcccccccc}
\toprule
 & \multicolumn{2}{c}{Res., random} & \multicolumn{2}{c}{Res., targeted} & \multicolumn{2}{c}{T$\to$T, random} & \multicolumn{2}{c}{T$\to$T, targeted} \\
\cmidrule(lr){2-3}\cmidrule(lr){4-5}\cmidrule(lr){6-7}\cmidrule(lr){8-9}
Score & in-situ & transf. & in-situ & transf. & in-situ & transf. & in-situ & transf. \\
\midrule
CoMention & 0.092 & 0.090 & 0.090 & 0.084 & \textbf{0.122} & \textbf{0.122} & \textbf{0.115} & \textbf{0.114} \\
BM25 & 0.101 & 0.098 & 0.101 & 0.095 & \textbf{0.119} & \textbf{0.110} & \textbf{0.119} & \textbf{0.112} \\
Dense & 0.100 & 0.097 & 0.100 & 0.096 & \textbf{0.116} & \textbf{0.111} & \textbf{0.116} & \textbf{0.111} \\
BGE dense & 0.101 & 0.098 & 0.101 & 0.094 & \textbf{0.106} & 0.093 & \textbf{0.106} & 0.093 \\
BGE rerank & 0.101 & 0.098 & 0.100 & 0.093 & \textbf{0.108} & 0.088 & \textbf{0.108} & 0.088 \\
LLM judge & 0.101 & 0.099 & 0.100 & 0.103 & 0.080 & 0.076 & 0.080 & 0.077 \\
NLI-verify & 0.101 & 0.099 & 0.100 & 0.099 & 0.101 & 0.101 & \textbf{0.110} & \textbf{0.107} \\
ESV-Scope & 0.099 & 0.094 & 0.099 & 0.094 & \textbf{0.117} & 0.104 & \textbf{0.115} & \textbf{0.108} \\
\midrule
CoMention +OA & 0.093 & 0.084 & 0.092 & 0.082 & \textbf{0.112} & 0.103 & \textbf{0.109} & 0.100 \\
Dense +OA & 0.101 & 0.097 & 0.101 & 0.094 & \textbf{0.111} & \textbf{0.107} & \textbf{0.111} & \textbf{0.107} \\
BGE rerank +OA & 0.100 & 0.099 & 0.100 & 0.091 & \textbf{0.125} & \textbf{0.121} & \textbf{0.125} & \textbf{0.111} \\
ESV-Scope +OA & 0.100 & 0.100 & 0.100 & 0.104 & \textbf{0.126} & \textbf{0.126} & \textbf{0.122} & \textbf{0.125} \\
\bottomrule
\end{tabular}

\end{table}

\section{Second domain: GapClose-ClimateFEVER}\label{sec:climate}

Table~\ref{tab:climate} gives the per-rate results behind the paper's Table~VII (resampled splits, random deletion). ESV-Scope has power zero at every rate because its hard numeric check sets the score of 19.5\% of \closed{} questions to zero, which pushes the conformal threshold to zero; its false novelty is then trivially zero as well.

\begin{table}[h]
\centering
\caption{GapClose-ClimateFEVER, resampled splits, random deletion, by deletion rate.}
\label{tab:climate}
\setlength{\tabcolsep}{2.5pt}
\begin{tabular}{lccccccc}
\toprule
Score & $0$ & $0.1$ & $0.2$ & $0.3$ & $0.5$ & $0.7$ & $1$ \\
\midrule
\multicolumn{8}{l}{\emph{False novelty, F1-optimal threshold}} \\
CoMention & \textbf{0.321} & \textbf{0.328} & \textbf{0.345} & \textbf{0.348} & \textbf{0.360} & \textbf{0.378} & \textbf{0.417} \\
BM25 & \textbf{0.383} & \textbf{0.384} & \textbf{0.396} & \textbf{0.409} & \textbf{0.428} & \textbf{0.468} & \textbf{0.538} \\
Dense & \textbf{0.294} & \textbf{0.299} & \textbf{0.314} & \textbf{0.319} & \textbf{0.346} & \textbf{0.366} & \textbf{0.432} \\
NLI-verify & \textbf{0.420} & \textbf{0.433} & \textbf{0.446} & \textbf{0.459} & \textbf{0.488} & \textbf{0.523} & \textbf{0.603} \\
ESV-Scope & \textbf{0.333} & \textbf{0.339} & \textbf{0.346} & \textbf{0.356} & \textbf{0.379} & \textbf{0.405} & \textbf{0.463} \\
\multicolumn{8}{l}{\emph{False novelty, Split conformal}} \\
CoMention & 0.094 & 0.095 & 0.098 & 0.099 & \textbf{0.107} & \textbf{0.114} & \textbf{0.129} \\
BM25 & 0.098 & 0.103 & \textbf{0.107} & \textbf{0.113} & \textbf{0.123} & \textbf{0.138} & \textbf{0.165} \\
Dense & 0.100 & 0.102 & 0.103 & \textbf{0.107} & \textbf{0.115} & \textbf{0.127} & \textbf{0.147} \\
NLI-verify & 0.099 & 0.103 & \textbf{0.110} & \textbf{0.118} & \textbf{0.138} & \textbf{0.163} & \textbf{0.240} \\
ESV-Scope & 0.000 & 0.000 & 0.000 & 0.000 & 0.000 & 0.000 & 0.000 \\
\multicolumn{8}{l}{\emph{False novelty, In-situ (Thm.~1)}} \\
CoMention & 0.094 & 0.094 & 0.093 & 0.092 & 0.092 & 0.092 & 0.092 \\
BM25 & 0.098 & 0.098 & 0.098 & 0.100 & 0.097 & 0.098 & 0.098 \\
Dense & 0.100 & 0.099 & 0.098 & 0.097 & 0.098 & 0.099 & 0.098 \\
NLI-verify & 0.099 & 0.097 & 0.098 & 0.098 & 0.099 & 0.099 & 0.099 \\
ESV-Scope & 0.000 & 0.000 & 0.000 & 0.000 & 0.000 & 0.000 & 0.000 \\
\multicolumn{8}{l}{\emph{False novelty, Transfer, $\hat p$ (Thm.~2)}} \\
CoMention & 0.094 & 0.092 & 0.092 & 0.090 & 0.090 & 0.090 & 0.092 \\
BM25 & 0.096 & 0.096 & 0.095 & 0.096 & 0.093 & 0.095 & 0.098 \\
Dense & 0.099 & 0.098 & 0.095 & 0.095 & 0.094 & 0.095 & 0.098 \\
NLI-verify & 0.097 & 0.093 & 0.093 & 0.093 & 0.092 & 0.089 & 0.099 \\
ESV-Scope & 0.000 & 0.000 & 0.000 & 0.000 & 0.000 & 0.000 & 0.000 \\
\multicolumn{8}{l}{\emph{Power, transfer calibration}} \\
CoMention & 0.145 & 0.143 & 0.140 & 0.137 & 0.133 & 0.128 & 0.116 \\
BM25 & 0.143 & 0.137 & 0.131 & 0.128 & 0.117 & 0.107 & 0.090 \\
Dense & 0.256 & 0.247 & 0.238 & 0.233 & 0.223 & 0.213 & 0.195 \\
NLI-verify & 0.107 & 0.099 & 0.095 & 0.091 & 0.077 & 0.060 & 0.021 \\
ESV-Scope & 0.000 & 0.000 & 0.000 & 0.000 & 0.000 & 0.000 & 0.000 \\
\bottomrule
\end{tabular}

\end{table}

\clearpage
\section{Retrospective validation}\label{sec:retro}

Table~\ref{tab:retro} reports all three cutoffs. ``Filled'' is the share of later-filled gaps certified \open{}, ``\textsc{nei}'' the share of \textsc{nei} questions certified \open{}, and ``Prec.'' the share of declared gaps that were indeed open at the cutoff. Table~\ref{tab:dossier} lists ten of the \SuppNdos{} gaps certified by ESV-Scope +OA at the 2010 cutoff in one run (threshold $\tau=\SuppTauDos{}$), chosen as those whose closest visible local evidence had the highest score. The lag between the cutoff and the answering paper ranges from \SuppLagMin{} to \SuppLagMax{} years among all \SuppNdos{} dossiers. The closest evidence is usually on a neighboring topic and scores far below the threshold, which illustrates why such questions look open at the cutoff.

\begin{table}[h]
\centering
\caption{Retrospective certification at three cutoffs, in-situ calibration, mean over \SuppR{} splits.}
\label{tab:retro}
\setlength{\tabcolsep}{3pt}
\begin{tabular}{lccccc}
\toprule
Score & F1-opt.\ FN & In-situ FN & Filled & \textsc{nei} & Prec. \\
\midrule
\multicolumn{6}{l}{\textit{Cutoff 2008: 45\% of documents visible, 221 closed, 415 later filled, 416 \textup{\textsc{nei}}}} \\
Dense & \textbf{0.378} & 0.101 & 0.499 & 0.501 & 0.966 \\
BGE rerank & \textbf{0.297} & 0.095 & 0.717 & 0.656 & 0.977 \\
LLM judge & \textbf{0.467} & 0.100 & 0.380 & 0.291 & 0.954 \\
ESV-Scope & \textbf{0.405} & 0.098 & 0.431 & 0.418 & 0.962 \\
BGE rerank +OA & \textbf{0.326} & 0.097 & 0.688 & 0.589 & 0.975 \\
ESV-Scope +OA & \textbf{0.394} & 0.102 & 0.420 & 0.392 & 0.959 \\
\multicolumn{6}{l}{\textit{Cutoff 2010: 56\% of documents visible, 310 closed, 326 later filled, 416 \textup{\textsc{nei}}}} \\
Dense & \textbf{0.340} & 0.095 & 0.426 & 0.426 & 0.941 \\
BGE rerank & \textbf{0.278} & 0.097 & 0.684 & 0.645 & 0.961 \\
LLM judge & \textbf{0.399} & 0.096 & 0.355 & 0.261 & 0.922 \\
ESV-Scope & \textbf{0.375} & 0.097 & 0.437 & 0.414 & 0.940 \\
BGE rerank +OA & \textbf{0.302} & 0.098 & 0.681 & 0.571 & 0.958 \\
ESV-Scope +OA & \textbf{0.340} & 0.099 & 0.426 & 0.358 & 0.935 \\
\multicolumn{6}{l}{\textit{Cutoff 2012: 68\% of documents visible, 407 closed, 229 later filled, 416 \textup{\textsc{nei}}}} \\
Dense & \textbf{0.352} & 0.100 & 0.375 & 0.384 & 0.893 \\
BGE rerank & \textbf{0.273} & 0.100 & 0.604 & 0.574 & 0.928 \\
LLM judge & \textbf{0.341} & 0.098 & 0.301 & 0.236 & 0.858 \\
ESV-Scope & \textbf{0.371} & 0.097 & 0.412 & 0.361 & 0.897 \\
BGE rerank +OA & \textbf{0.325} & 0.097 & 0.592 & 0.497 & 0.924 \\
ESV-Scope +OA & \textbf{0.326} & 0.097 & 0.401 & 0.306 & 0.889 \\
\bottomrule
\end{tabular}

\end{table}

\begin{table}[h]
\centering
\footnotesize
\caption{Gap dossiers at the 2010 cutoff (ESV-Scope +OA, in-situ calibration). Lag in years from the cutoff.}
\label{tab:dossier}
\begin{tabular}{>{\raggedright\arraybackslash}p{3.0cm}>{\raggedright\arraybackslash}p{4.7cm}>{\raggedright\arraybackslash}p{6.6cm}c}
\toprule
Certified open question & Closest local evidence (year, score) & Later answering paper & Lag \\
\midrule
Trans-acting factors, such as lncRNAs, influence mRNA translation. & ``We now show that HCV infection triggers phosphorylation and activation of the RNA-dependent protein kinase PKR, which \ldots{}'' (2009, 0.028) & ``Here we report a posttranscriptional function for lincRNA-p21 as a modulator of translation.'' (support, 2012, doi:10.1016/j.molcel.2012.06.027) & 2 \\
HbA1c consistently overestimates glycemic control in black individuals with Sickle cell trait (SCT). & ``In seven studies reporting early gestational glycosylated hemoglobin values, mean levels were consistently lower in PCC \ldots{}'' (2002, 0.026) & ``In unadjusted GEE analyses, for a given fasting glucose, HbA1c values were statistically significantly lower in those with (5.72\%) vs those without (6.01\%) \ldots{}'' (contradict, 2017, doi:10.1001/jama.2016.21035) & 7 \\
Ambulatory blood pressure monitoring is inaccurate at diagnosing hypertension. & ``The CV value of daytime SBP and the SD value of both 24-h SBP and daytime SBP also had positive correlations with the \ldots{}'' (2005, 0.025) & ``Compared with ambulatory monitoring thresholds of 135/85 mm Hg, clinic measurements over 140/90 mm Hg had mean sensitivity and specificity of 74.6\% (95\% \ldots{}'' (contradict, 2011, doi:10.1136/bmj.d3621) & 1 \\
AMP-activated protein kinase (AMPK) activation reduces inflammation-related fibrosis in the lungs. & ``In isolated rat skeletal muscles, metformin stimulates glucose uptake coincident with AMPK activation.'' (2001, 0.019) & ``Pharmacological activation of AMPK in myofibroblasts from lungs of humans with IPF display lower fibrotic activity, along with enhanced mitochondrial \ldots{}'' (support, 2018, doi:10.1038/s41591-018-0087-6) & 8 \\
Increased microtubule acetylation exacerbates LRRK2 Roc-COR domain mutation induced locomotor deficits. & ``In contrast, mutation in the P2 or P4 site did not affect the apparent binding affinity of the mutant for microtubules \ldots{}'' (2004, 0.019) & ``In vitro, increasing microtubule acetylation using deacetylase inhibitors or the tubulin acetylase  TAT1 prevents association of mutant LRRK2 with \ldots{}'' (contradict, 2014, doi:10.1038/ncomms6245) & 4 \\
Lice attenuated SIV vaccines induce a stronger antigen-specific T cell response in lymph node cells. & ``Using nanoparticles conjugated to the model antigen ovalbumin, we demonstrate generation of humoral and cellular \ldots{}'' (2007, 0.019) & ``Here we show that the degree of LAV-mediated protection against intravenous wild-type SIVmac239 challenge strongly correlates with the magnitude and function \ldots{}'' (support, 2012, doi:10.1038/nm.2934) & 2 \\
Auditory entrainment is strengthened when people see congruent visual and auditory information. & ``Blind individuals often demonstrate enhanced nonvisual perceptual abilities.'' (2005, 0.019) & ``UNLABELLED Congruent audiovisual speech enhances our ability to comprehend a speaker, even in noise-free conditions.'' (support, 2015, doi:10.1523/JNEUROSCI.1829-15.2015) & 5 \\
Ca2+ cycling is a UCP1-independent thermogenic mechanism. & ``In immune cells, Ca2+ signaling is an essential step in many forms of activation.'' (2009, 0.018) & ``Here we report a robust UCP1-independent thermogenic mechanism in beige fat that involves enhanced ATP-dependent Ca2+ cycling by sarco/endoplasmic reticulum \ldots{}'' (support, 2017, doi:10.1038/nm.4429) & 7 \\
In British Men, haplogroup I increases risk of cardiovascular disease by 50\%. & ``RESULTS Women were older than men and had higher rates of cardiovascular risk factors.'' (2008, 0.018) & ``Carriers of haplogroup I had about a 50\% higher age-adjusted risk of coronary artery disease than did men with other Y chromosome lineages in BHF-FHS (odds \ldots{}'' (support, 2012, doi:10.1016/S0140-6736(11)61453-0) & 2 \\
Reducing H3k4me3 methylation induces mouse epiblast stem cells to naive pluripotency efficiently. & ``Finally, the same growth conditions allow the derivation of human induced pluripotent stem (iPS) cells with similar \ldots{}'' (2010, 0.018) & ``Here, we show that blocking histone H3K4 methyltransferase MLL1 activity with the small-molecule inhibitor MM-401 reprograms mouse epiblast stem cells (EpiSCs) \ldots{}'' (support, 2016, doi:10.1016/j.stem.2016.02.004) & 6 \\
\bottomrule
\end{tabular}

\end{table}

\section{Reproduction}
The steps below reproduce every number in the paper and in this document; they are also listed in the repository README.
\begin{verbatim}
cd ESV-Gap/gapclose/src
python build_benchmark.py
python score_candidates.py --tag _large --batch 32 \
    --nli MoritzLaurer/DeBERTa-v3-large-mnli-fever-anli-ling-wanli
python cv_dev.py --tag _large && python evaluate.py --tag _large
python bge_scores.py && python llm_judge.py && python extra_uncertified.py
python fetch_external.py             # needs OPENALEX_API_KEY
python certify.py --protocol pooled  --deletion mcar   # and targeted
python certify.py --protocol holdout --deletion mcar   # and targeted
python fetch_years.py && python retro.py
# second domain: same scripts with GAPCLOSE_OUT=../outputs_climate
GAPCLOSE_OUT=../outputs_climate python build_climate.py
cd ../../paper_gapclose_icai2026
python make_tables.py && python make_supp.py
pdflatex main.tex && pdflatex supplementary.tex
\end{verbatim}

\end{document}
